\documentclass[11pt]{article}
\usepackage{amsmath,amssymb,amsthm}
\usepackage[margin=1in]{geometry}
\usepackage{hyperref}
\usepackage{listings}
\usepackage{xcolor}

\lstset{
  basicstyle=\ttfamily\small,
  keywordstyle=\color{blue},
  commentstyle=\color{gray},
  showstringspaces=false,
  breaklines=true,
}

\newtheorem{theorem}{Theorem}
\newtheorem{lemma}[theorem]{Lemma}
\newtheorem{definition}[theorem]{Definition}
\newtheorem{remark}[theorem]{Remark}
\newtheorem{proposition}[theorem]{Proposition}

\title{Formalization of G\"odel's Incompleteness Theorems\\
  in Basic Recursive Arithmetic (BRA)\\
  \large{BRA2 collapsed-Tree edition}}
\author{Summary of an Agda development}
\date{May 2026}

\begin{document}
\maketitle

\begin{abstract}
We present a fully formalized proof of both G\"odel's First and Second
Incompleteness Theorems in a tree-based variant of Guard's Basic
Recursive Arithmetic (BRA), executed in Agda~2.8 with
\verb|--safe --without-K --exact-split|.  The development contains
\textbf{zero postulates and zero holes}; \verb|godelII : Deriv Con -> Deriv bot|
typechecks end-to-end against the abstract Goedel~II skeleton.

This is the BRA2 edition.  Compared with the earlier BRA development,
the Layer-1/Layer-2 split has been collapsed: there is now a single
\verb|Term| datatype (no separate \verb|Tree| datatype), and ``codes''
are exactly the value-shaped Terms certified by an
\verb|IsValue : Term -> Set| predicate (built only from $O$ and
\verb|ap2 Pair|).  The map \verb|reify : Tree -> Term| of the prior
edition becomes the identity; every place the old code wrote
\verb|reify y| is now just \verb|y|.  The substantive math is unchanged;
the bookkeeping is uniform.

The headline results are:
\begin{itemize}
  \item $\mathit{thm11} : \mathit{Deriv}\;G \to \mathit{Deriv}\;\bot$
        (G\"odel~I);
  \item $\mathit{godelII} : \mathit{Deriv}\;\mathit{Con} \to \mathit{Deriv}\;\bot$
        (G\"odel~II).
\end{itemize}
The development follows Guard's 1963 lecture
notes~\cite{guard1963}.  The intermediate construction is a single
first-order combinator $\mathit{constr5} : \mathit{Fun}_1$ realising
Guard's ``$F(x)$'' (page~17); each step is given by an explicit
equational derivation.

A key conceptual point this note emphasises (Section~\ref{sec:internal})
is that intermediate steps of Theorem~14's chain produce
\emph{equational} \verb|Deriv|s of the form
\verb|atomic (eqn (thmT t) u)| where \verb|u| is \emph{not}
\verb|codeFormula P| for any $P : \mathit{Formula}$.  Only the final
\verb|step5| conclusion has \verb|u = codeFormula bot| (with $\bot$
genuinely closed).  Guard's chain operates one level lower than
Hilbert--Bernays--L\"ob: instead of internal provability claims about
specific formulas, it manipulates substituted-codes through standard
BRA inference rules.  This makes the entire proof internal and
equational --- a feature of the framework, not a gap.
\end{abstract}

\section{Introduction}

The aim is a feasible, ASCII-only, syntax-first Agda formalization of
G\"odel's two incompleteness theorems following Guard's 1963 system of
\emph{Basic Recursive Arithmetic}~\cite{guard1963}.

\subsection{Single-layer Term language with $\mathit{IsValue}$}
\label{subsec:layers}

In BRA2 there is one syntactic layer.

\paragraph{The \texttt{Term} language.}
\begin{itemize}
  \item $\mathit{Term}$: built from $O$ (leaf), $\mathit{var}\;n$,
        $\mathit{ap}_1\;f\;t$, $\mathit{ap}_2\;g\;t_1\;t_2$.
  \item Function symbols: $\mathit{Fun}_1$
        ($\{I, \mathit{Fst}, \mathit{Snd}, \mathit{Comp}, \mathit{Comp}_2, Z\}$,
        plus $\mathit{KT}\,\_$ as a defined function), $\mathit{Fun}_2$
        ($\{\mathit{Pair}, \mathit{Const}, \mathit{Lift}, \mathit{Post}, \mathit{Fan}, \mathit{IfLf}, \mathit{TreeEq}, \mathit{treeRec}\}$).
  \item $\mathit{Formula}$: $\mathit{atomic}\;(\mathit{eqn}\;\ell\;r)$, $\mathit{not}\;\_$, $\_\,\mathit{imp}\,\_$.
  \item $\mathit{Deriv}\;F$: inductive family of derivations.  Constructors
        for the algebraic axioms ($\mathit{ax}_*$), the equality axioms
        (Guard's Ax~4--7), three propositional axioms ($\mathit{axK}, \mathit{axS}, \mathit{axNeg}$
        in Lukasiewicz form), modus ponens ($\mathit{mp}$), substitution
        ($\mathit{ruleInst}$), and binary tree induction ($\mathit{ruleIndBT}$).
\end{itemize}

\paragraph{Codes are value-shaped Terms.}
``Codes'' (formerly Layer~2 \verb|Tree|s) are now Terms whose value-shape
is certified by
\[
  \mathit{IsValue} : \mathit{Term} \to \mathit{Set},
\]
with constructors
$\mathit{valO} : \mathit{IsValue}\;O$ and
$\mathit{valNd}\;a\;b\;v_a\;v_b : \mathit{IsValue}\;(\mathit{ap}_2\;\mathit{Pair}\;a\;b)$
(when $\mathit{IsValue}\;a$ and $\mathit{IsValue}\;b$).

The encoding maps now have type
\[
  \mathit{code} : \mathit{Term} \to \mathit{Term}, \qquad
  \mathit{codeFormula} : \mathit{Formula} \to \mathit{Term},
\]
together with structural-soundness lemmas
$\mathit{code\_isValue} : (t : \mathit{Term}) \to \mathit{IsValue}\;(\mathit{code}\;t)$,
$\mathit{codeFormula\_isValue}$, $\mathit{codeF}_1\_\mathit{isValue}$,
$\mathit{codeF}_2\_\mathit{isValue}$, and $\mathit{natCode\_isValue}$
(BRA2/Term.agda, BRA2/Formula.agda).  Every primitive constructor of an
encoding produces a value-shaped Term.

\paragraph{The legacy \texttt{Tree}, \texttt{lf}, \texttt{nd}, \texttt{reify}.}
For continuity with Guard's prose and with the BRA edition,
BRA2/Term.agda exposes
\[
  \mathit{Tree} \;:=\; \mathit{Term}, \quad
  \mathit{lf} \;:=\; O, \quad
  \mathit{nd}\;a\;b \;:=\; \mathit{ap}_2\;\mathit{Pair}\;a\;b, \quad
  \mathit{reify}\;t \;:=\; t.
\]
Old code that wrote ``\verb|reify y|'' for $y : \mathrm{Tree}$ now
literally evaluates to $y$.  Conceptually the Layer-1/Layer-2 split is
preserved as the $\mathit{IsValue}$ predicate; syntactically the layers
are unified.

\paragraph{Reflection layer (object-level operations on codes).}
Functors that reason about BRA syntax internally:
\begin{itemize}
  \item $\mathit{cor} : \mathit{Fun}_1$: defined as $\mathit{Rec}\;\mathit{stepCor}$
        (BRA2/Cor.agda).  Spec:
        $\mathit{corOfReify} : (t : \mathit{Term}) \to \mathit{IsValue}\;t \to
          \mathit{Deriv}\;(\mathit{ap}_1\;\mathit{cor}\;t = \mathit{code}\;t)$.
        Operationally, $\mathit{cor}$ takes a value-shaped Term (= a
        ``numeral'') and returns its code.
  \item $\mathit{sub} : \mathit{Fun}_2$, $\mathit{subT} : \mathit{Fun}_2$:
        coded substitution (BRA2/Sub.agda, BRA2/SubT.agda), now indexed
        by $\mathit{IsValue}$ on the substituent.
  \item $\mathit{thmT} : \mathit{Fun}_1$: the proof-checker
        (Section~\ref{sec:thmT}).
\end{itemize}

\paragraph{Reflection theorem.}
\[
  \mathit{encode} : (P : \mathit{Formula}) \to \mathit{Deriv}\;P \to
    \Sigma\;\mathit{Term}\;\bigl(\lambda y.\;
      \Sigma\;(\mathit{IsValue}\;y)\;\bigl(\lambda \_.\;\mathit{Deriv}\;
        (\mathit{thmT}\;y = \mathit{codeFormula}\;P)\bigr)\bigr).
\]
The new $\mathit{IsValue}\;y$ field is a structural strengthening over
the BRA edition: every encoded proof tree is now certified value-shaped,
so $\mathit{cor}\;y$ provably reduces to $\mathit{code}\;y$ at the witness
returned by $\mathit{encode}$.  This certification is exactly what
discharges the parametric \verb|cor x| in the chain at the closure
(see Section~\ref{sec:internal}).

\subsection{Notational convention: code-quotes}
\label{subsec:codequote}

We use $\ulcorner E \urcorner$ throughout for ``the code (value-shaped
Term) of $E$'':
\[
  \ulcorner t \urcorner := \mathit{code}\;t \quad (t : \mathit{Term}), \qquad
  \ulcorner F \urcorner := \mathit{codeFormula}\;F \quad (F : \mathit{Formula}).
\]
Whenever a Guard-style expression mixes object-level and encoded
material (e.g.\ $f(\underline{x})$), we keep the code-quote explicit:
the formal object is the Term $\ulcorner f(\underline{x}) =
\underline{f(x)}\urcorner$, never the unquoted equation.  See
Section~\ref{subsec:underline-caveat} for the full caveat.

\section{The provability predicate $\mathit{thmT}$}
\label{sec:thmT}

The G\"odel-style ``$\mathit{th}$'' provability predicate is realized
in BRA by a single $\mathit{Fun}_1$-functor:
\[
  \mathit{thmT} \;=\; \mathit{Rec}\;\mathit{stepProtoWrapped} \;:\; \mathit{Fun}_1.
\]
Here $\mathit{stepProtoWrapped}$ is a $\mathit{Fun}_2$ assembled from a
40-deep linear $\mathit{IfLf}$ cascade (BRA2/Thm/StepProto.agda and
BRA2/Thm/ThmT.agda), one branch per primitive proof rule: ten
equational axioms of Group~I, eight Group~II axioms, nine
equality and propositional rules, and three rules of inference
($\mathit{mp}$, $\mathit{ruleInst}$, $\mathit{ruleIndBT}$).

The \emph{key} property of $\mathit{thmT}$ is the encoding-soundness
lemma $\mathit{encode}$ above, proved by induction on the
$\mathit{Deriv}$ derivation in BRA2/Thm/ThmT.agda's $\mathit{encodeRich}$.

\subsection{Self-substitution closure}

$\mathit{thmT}$ is a closed combinator:
\[
  \mathit{thClosed} : \forall x\;r,\quad \mathit{substF}_1\;x\;r\;\mathit{thmT} = \mathit{thmT}
\]
and its encoding is var-free:
\[
  \mathit{codeF1Th\_noVar} : \forall x\;\mathit{repl},\quad
    \mathit{codedSubst}\;\mathit{repl}\;(\mathit{code}\;(\mathit{var}\;x))\;(\mathit{codeF}_1\;\mathit{thmT})
    = \mathit{codeF}_1\;\mathit{thmT}.
\]

\section{G\"odel's First Incompleteness Theorem (Theorem~11)}

\subsection{The diagonal sentence}

Following Guard, define
\begin{align*}
  F_{\mathrm{pre}} \;&:=\; \neg \big(\mathit{thmT}\;v_1 \;=\; \mathit{sub}\;v_0\;v_0\big),\\
  j \;&:=\; \mathit{codeFormula}\;F_{\mathrm{pre}} \quad \text{(a closed value-shaped Term)},\\
  G \;&:=\; F_{\mathrm{pre}}[v_0 := j],
\end{align*}
with $i := j$ (since $\mathit{reify} = \mathrm{id}$ in BRA2).  The
sentence $G$ has $v_1$ free; informally, $G$ asserts of itself that no
proof tree $v_1$ proves the formula whose code is $j[v_0 := j]$.

The new ingredient compared to BRA: $j\_\mathit{isValue} :
\mathit{IsValue}\;j$ is exposed inside the
\verb|Diagonal| abstract block (BRA2/Thm11Diagonal.agda).  This
certificate propagates through $\mathit{Outer.GodelII}$ and is what
makes \verb|corOfReify j j_isValue| usable downstream.

\subsection{Theorem 11}

\begin{theorem}[\texttt{thm11}, G\"odel I]
\label{thm:godel-I}
\(
  \mathit{thm11} : \mathit{Deriv}\;G \to \mathit{Deriv}\;\bot,
\)
where $\bot := \mathit{atomic}\;(\mathit{eqn}\;\mathit{trueT}\;\mathit{falseT})$,
$\mathit{trueT} := O$, $\mathit{falseT} := \mathit{ap}_2\;\mathit{Pair}\;O\;O$.
\end{theorem}

\begin{proof}[Sketch]
From any $\pi : \mathit{Deriv}\;G$ extract its encoding
$(y_d, v_d, E_1)$ where $v_d : \mathit{IsValue}\;y_d$ and
$E_1 : \mathit{Deriv}\;(\mathit{thmT}\;y_d = \mathit{codeFormula}\;G)$.
A separate diagonal lemma (BRA2/Thm11Diagonal.agda) yields
$N : \mathit{Deriv}\;\neg(\mathit{thmT}\;y_d = \mathit{codeFormula}\;G)$,
and $\mathit{ax}_{\mathit{ExFalso}}$ + two $\mathit{mp}$'s collapse $E_1$
and $N$ to $\mathit{Deriv}\;\bot$.
\end{proof}

\section{G\"odel's Second Incompleteness Theorem (Theorem~14)}

\begin{theorem}[\texttt{godelII}, G\"odel II]
\label{thm:godel-II}
\(
  \mathit{godelII} : \mathit{Deriv}\;\mathit{Con} \to \mathit{Deriv}\;\bot,
\)
where $\mathit{Con} := \neg(\mathit{thmT}\;v_0 = \mathit{codeFormula}\;\bot)$.
\end{theorem}

\subsection{Caveat: ``$f(\underline{x}) = \underline{f(x)}$'' is not a BRA formula}
\label{subsec:underline-caveat}

A structurally crucial point, easy to miss in Guard's notation: the
expression
\[
  f(\underline{x}) \;=\; \underline{f(x)}
\]
is \textbf{not} a BRA formula.  Guard's underline $\underline{(\cdot)}$
takes a value and returns its numeral term --- well-defined only at the
meta-level for closed numerical inputs.  For a variable $x$ (the case
Theorems~12--14 actually need), neither $\underline{x}$ nor
$\underline{f(x)}$ has any BRA referent: there is no operator in BRA's
term language that, given an arbitrary $\mathrm{Term}$, returns ``the
numeral of its value.''

The only formal object that \emph{does} correspond to Guard's
expression is its \emph{code}, the Term
\[
  \ulcorner f(\underline{x}) \;=\; \underline{f(x)} \urcorner,
\]
which Definition~12 of Guard's notes \emph{defines} by stipulation.  In
BRA2 this Term is built parametrically as $\mathit{codeFTeq}_1\;f\;x$,
with $\mathit{cor}\;x$ filling the slot for the code of $x$'s numeral
and $\mathit{cor}\;(f\,x)$ filling the slot for the code of $f(x)$'s
numeral.

When $x$ is a free Term variable, neither $\mathit{cor}\;x$ nor
$\mathit{cor}\;(f\,x)$ reduces (the defining equation
$\mathit{corOfReify}$ fires only on $\mathit{IsValue}$ inputs), so
$\mathit{codeFTeq}_1\;f\;x$ contains $\mathit{ap}_1$-headed subterms
literally.  Codes (in the sense of $\mathit{codeFormula}\;P$) are
$\mathit{IsValue}$-shape, so $\mathit{codeFTeq}_1\;f\;x$ at variable $x$
is \textbf{not} of the form $\mathit{codeFormula}\;P$ for any
$P : \mathit{Formula}$.

\subsection{Theorem 12 --- the encoded ``$\mathit{th}\;f$'' clause}

For every $\mathit{Fun}_1$ functor $f$, Theorem~12 produces a
$\mathit{Fun}_1$ $D_f$ such that $\mathit{thmT}$ evaluated at $D_f\;x$ is
provably equal to the encoded form of
``$f(\underline{x}) = \underline{f(x)}$'':

\begin{theorem}[$\mathit{thm12}\_F_1$]
\(
  \mathit{thm12}\_F_1 : (f : \mathit{Fun}_1) \to \mathit{Thm12\_F_1\_Result}\;f
\)
where $\mathit{Thm12\_F_1\_Result}\;f$ packages $D_f : \mathit{Fun}_1$ with
\(
  \mathit{proof} : (x : \mathit{Term}) \to
   \mathit{Deriv}\;\bigl(\mathit{thmT}\;(D_f\;x)
       \;=\; \mathit{codeFTeq}_1\;f\;x\bigr).
\)
\end{theorem}
The conclusion's RHS is the literal Pair-shape encoded equation.  An
analogous statement for $\mathit{Fun}_2$.  Structural recursion on $f$
with one clause per constructor (15 total),
BRA2/Thm12.agda module $\mathit{FromBridges}$.

\subsection{Theorem 13 --- under hypothesis $f\;x = y$}

\begin{lemma}[$\mathit{thm13}\_F_1$]
For every $f : \mathit{Fun}_1$ and every Theorem~12 result $r_{12}$ for $f$,
\(
  \mathit{thm13}\_F_1\;f\;r_{12}\;x\;y : \mathit{Deriv}\;(f\;x = y) \to
    \mathit{Deriv}\;(\mathit{thmT}\;(D_f\;x) = \mathit{codeFTeq}_{1,\mathrm{Hyp}}\;f\;x\;y).
\)
\end{lemma}
The conclusion replaces $\mathit{cor}\;(f\,x)$ by $\mathit{cor}\;y$.
One-line proof: $\mathit{ruleTrans}$ on $r_{12}.\mathit{proof}\;x$ with
$\mathit{cong}_R\;\mathit{Pair}$ on $\mathit{cong}_1\;\mathit{cor}$ of
the hypothesis.

\subsection{Theorem 14 --- the consistency-to-bot construction}

The contract abstracted in BRA2/Thm14Abstract.agda asks for:
\begin{enumerate}
  \item a $\mathit{Fun}_1$-functor $\mathit{constr5}$ such that
        $\mathit{thmT}\;(\mathit{constr5}\;x)$ syntactically encodes
        $\bot$;
  \item a parametric internal-implication
\[
  \mathit{step5} : (x : \mathit{Term}) \to
   \mathit{Deriv}\bigl((\mathit{thmT}\;x = \mathit{codeFormula}\;G)
                  \;\mathit{imp}\;
                  (\mathit{thmT}\;(\mathit{constr5}\;x) = \mathit{codeFormula}\;\bot)\bigr).
\]
\end{enumerate}

The closure $\mathit{closureToG}$ instantiates $\mathit{step5}$ at $v_1$,
$\mathit{ruleInst}$s $\mathit{Con}$, modus-ponenses, transports through
$\mathit{subIIeqcG}$ and $G_{\mathrm{unfold}}$, arriving at
$\mathit{Deriv}\;G$.  Composing with $\mathit{thm11}$ gives $\mathit{godelII}$.

\subsubsection{Construction of $\mathit{constr5}$}

Two pre-existing tautologies, encoded:
\begin{itemize}
  \item $t := (v_0 = v_2)\;\mathit{imp}\;((v_1 = v_2)\;\mathit{imp}\;(v_0 = v_1))$ (transitivity);
  \item $t' := (v_0 = v_1)\;\mathit{imp}\;((\neg(v_0 = v_1))\;\mathit{imp}\;\bot)$ (ex falso).
\end{itemize}
Both BRA-derivable (BRA2/Thm14Numerals.agda); $\mathit{encode}$ yields
value-shaped Terms $\mathit{tterm}, \mathit{t'term}$ certified by
$t\_term\_isValue$ / $t'\_term\_isValue$ (these certificates are
load-bearing for the chain --- see Section~\ref{sec:internal}).

The chain, paralleling Guard's page~17:
\begin{align*}
  f\_\mathit{part}(x) \;&:=\; \mathit{ruleInst}\;v_0\;(\mathit{th}\;\hat x)\;\big(\mathit{ruleInst}\;v_1\;(\mathit{sub}\;i\;i)\;(\mathit{ruleInst}\;v_2\;\mathit{cor}\;G\;\mathit{tterm})\big),\\
  g\_\mathit{part}(x) \;&:=\; \mathit{mp}\big(\mathit{mp}(f\_\mathit{part}(x), D_{\mathit{thmT}}\;x), D_{\mathit{sub}}\;i\;i\big),\\
  K\_\mathit{part}(x) \;&:=\; \mathit{ruleInst}\;v_1\;(\mathit{cor}\;x)\;x,\\
  h\_\mathit{part}\_\mathit{skel}(x) \;&:=\; \mathit{ruleInst}\;v_0\;(\mathit{th}\;\hat x)\;\big(\mathit{ruleInst}\;v_1\;(\mathit{sub}\;i\;i)\;\mathit{t'term}\big),\\
  \mathit{constr5}_{\mathrm{final}}(x) \;&:=\; \mathit{mp}\big(\mathit{mp}(h\_\mathit{part}\_\mathit{skel}(x), g\_\mathit{part}(x)), K\_\mathit{part}(x)\big).
\end{align*}

\subsubsection{The five steps}

(BRA2/Th14*.agda.)
\begin{description}
  \item[Step 1 --- $\mathit{step1\_l}$] Theorem~13 applied to $\mathit{thmT}$.
  \item[Step 2 --- $\mathit{step2\_l}$] Theorem~13 applied to $\mathit{sub}$ at $(i, i, \mathit{cor}\;G)$ under closed hypothesis $\mathit{subIIeqcG}$.
  \item[Step 3 --- $\mathit{step3\_l}$] Three sb-applications on $t$, canonicalised, then two $\mathit{mp}$ dispatches.
  \item[Step 4 --- $K\_\mathit{part}\_l$] One sb-application at $v_1$ on $G_{\mathrm{norm}}$ via the $\mathit{PrAtX}\;x$ hypothesis to swap $\mathit{thmT}\;x \mapsto \mathit{cor}\;G$.
  \item[Step 5 --- $\mathit{step5\_l}$] Two outer $\mathit{mp}$ dispatches at $\mathit{constr5}_{\mathrm{inner,final}}$ and $\mathit{constr5}_{\mathrm{final}}$.
\end{description}

\section{What is going on at the encoded layer: a remarkable internal proof}
\label{sec:internal}

\subsection{The phenomenon}

The chain $\mathit{step1\_l},\dots,\mathit{step5\_l}$ proves
$\mathit{Deriv}$s of the schematic form
\[
  \mathit{Deriv}\;\bigl(\mathit{PrAtX}\;x \;\mathit{imp}\;
    \mathit{atomic}\;(\mathit{eqn}\;(\mathit{thmT}\;t_i)\;u_i)\bigr),
\]
where $t_i, u_i$ are Terms depending on $x$.  Several of these
intermediate $u_i$ are \textbf{not} of the form $\mathit{codeFormula}\;P$
for any $P : \mathit{Formula}$.

A representative example: $\mathit{step3\_l}$ (BRA2/Th14Step3.agda)
concludes
\[
  \mathit{Deriv}\;\bigl(\mathit{PrAtX}\;x \;\mathit{imp}\;
    \mathit{atomic}\;(\mathit{eqn}\;(\mathit{thmT}\;(g\_\mathit{part}\;x))\;
      \mathit{ap}_2\;\mathit{Pair}\;(\widehat{\mathit{th}\;\hat x})\;\widehat{\mathit{sub}\;i\;i})\bigr),
\]
with
\[
  \widehat{\mathit{th}\;\hat x}
   \;=\; \mathit{ap}_2\;\mathit{Pair}\;\mathit{tagAp}_1\;
        (\mathit{ap}_2\;\mathit{Pair}\;(\mathit{codeF}_1\;\mathit{thmT})\;(\mathit{ap}_1\;\mathit{cor}\;x)).
\]
At a variable $x$ (e.g.\ $x = v_1$), $\mathit{ap}_1\;\mathit{cor}\;x$
does not reduce: $\mathit{corOfReify}$ requires $\mathit{IsValue}\;x$,
and $v_1$ is not value-shape.  The Term
$\mathit{ap}_1\;\mathit{cor}\;x$ has an $\mathit{ap}_1$ head; codes
are exclusively $\mathit{IsValue}$-shape (built from $O$ and
$\mathit{ap}_2\;\mathit{Pair}$).  Therefore $u_3 = \mathit{ap}_2\;
\mathit{Pair}\;(\widehat{\mathit{th}\;\hat x})\;\widehat{\mathit{sub}\;i\;i}$
is not in the image of $\mathit{codeFormula}$.

The same phenomenon recurs in $\mathit{step4}, \mathit{step5}_{\mathrm{intermediate}},
h\_\mathit{part}\_\mathit{skel}\_l$.  In each case, $u_i$ is a
\emph{substituted-code} --- the result of inserting $\mathit{cor}\;x$
into a variable slot of the code of an open formula.

Yet $\mathit{Deriv}\;(\mathit{atomic}\;(\mathit{eqn}\;t\;u))$ is
perfectly well-formed for arbitrary Terms $t, u$ --- the equational
predicate $\mathit{eqn}$ does not require $u$ to be a code of a formula.
The chain composes via the standard inference rules ($\mathit{mp}$,
$\mathit{ruleTrans}$, $\mathit{cong}_*$, $\mathit{axS}$, $\mathit{axK}$)
without ever needing the intermediate $u_i$ to be $\mathit{codeFormula}\;P$.

\subsection{Why this works: substituted-codes vs.\ codes-of-formulas}

There are two distinct readings of Guard's notation
$\ulcorner f(\underline{x}) \;=\; \underline{f(x)} \urcorner$:

\begin{description}
  \item[(A) Code-of-substituted-formula.]  Read $\underline{x}$ as
    $\mathit{cor}\;x$ (a BRA Term), and form the BRA formula
    $\mathit{atomic}\;(\mathit{eqn}\;(\mathit{ap}_1\;f\;(\mathit{cor}\;x))\;\dots)$.
    Take its $\mathit{codeFormula}$.
    The inner slot becomes $\mathit{code}\;(\mathit{cor}\;x)$
    (= the code of the BRA term $\mathit{cor}\;x$).
  \item[(B) Substituted-code.]  Take the open formula
    $\mathit{atomic}\;(\mathit{eqn}\;(\mathit{ap}_1\;f\;v_n)\;\dots)$,
    code it, then $\mathit{sb}$-substitute $\mathit{cor}\;x$ into the
    $\ulcorner v_n\urcorner$ slot.  The inner slot becomes
    $\mathit{cor}\;x$ itself.
\end{description}
At a free variable $x$ these readings disagree: (A) places
$\mathit{code}\;(\mathit{cor}\;x)$ in the slot, (B) places
$\mathit{cor}\;x$.  Theorem~12's $\mathit{codeFTeq}_1$ uses reading
(B) --- the substituted-code form.  This is forced by the structure of
$D_f$, which is built from $\mathit{Comp}_2\;\mathit{Pair}\;\dots\;\mathit{cor}$,
producing $\mathit{cor}\;x$ in its evaluation directly.

The two readings coincide \textbf{when $x$ is value-shaped}:
$\mathit{corOfReify}$ then gives
$\mathit{cor}\;x = \mathit{code}\;x$, and \emph{Lemma~codeCommutesFormula}
(BRA2/CodeCommutes.agda) gives
\[
  \mathit{code}\;(\mathit{substF}\;n\;t\;P) \;=\;
   \mathit{codedSubst}\;(\mathit{code}\;t)\;(\mathit{code}\;(\mathit{var}\;n))\;(\mathit{codeFormula}\;P).
\]
So at value-shape $x$, the substituted-code (B) is BRA-equal to
$\mathit{codeFormula}\;(\mathit{substF}\;n\;x\;P_{\mathrm{open}})$ ---
a genuine code-of-formula.  At variable $x$ this collapse fails, and
(B) remains a hybrid Term.

\subsection{The chain's discharge mechanism}

Why does Goedel~II close even though step~3 etc.\ produce hybrid RHS?

\begin{enumerate}
  \item \emph{Internal composition stays at the equational level.}
        The chain's mp/cong/ruleTrans steps operate on $\mathit{eqn}\;t\;u$
        for arbitrary $t, u$; they do not require $u$ to be
        $\mathit{codeFormula}\;P$.  All intermediate Derivs are honest
        BRA derivations producing well-typed atomic equations.
  \item \emph{$\mathit{step5}$'s conclusion is genuinely
        $\mathit{codeFormula}\;\bot$.}  At step~5 the final RHS is
        $\mathit{codeFormula}\;\bot$, where $\bot$ is closed.  No
        $\mathit{cor}\;x$ remains in the conclusion --- the
        Comp2-of-KT-of-codeFormula-bot construction in
        $\mathit{constr5}_{\mathrm{final}}$ is a constant in $x$
        modulo Lift, so its $\mathit{thmT}$-image is a closed code by
        construction.
  \item \emph{Closure transports through $\mathit{notEqTransport}$ at
        a value-shape witness.}  $\mathit{closureToG}$ instantiates
        $\mathit{step5}$ via $\mathit{ruleInst}$ at the encode-witness
        of $G$'s proof.  By the strengthened $\mathit{encode}$ contract
        (returning $\mathit{IsValue}$), this witness is a value-shaped
        Term.  At such a witness, $\mathit{cor}\;x$ inside any
        intermediate substituted-code rewrites to $\mathit{code}\;x$ via
        $\mathit{corOfReify}$, and the substituted-codes collapse to
        genuine $\mathit{codeFormula}$s of substituted formulas.
\end{enumerate}

The interleaving --- carry hybrid forms through internal-implication
chains, collapse at the closure --- is what lets the proof be
parametric over $(x : \mathit{Term})$ rather than meta-iterated over
closed numerals.

\subsection{Comparison with Hilbert--Bernays--L\"ob}

Standard textbook G\"odel~II proofs go through provability statements:
\[
  \text{D1} : \vdash A \;\Longrightarrow\; \vdash \mathrm{Pr}(\ulcorner A\urcorner),
  \quad
  \text{D2} : \vdash \mathrm{Pr}(\ulcorner A\to B\urcorner)
              \to \mathrm{Pr}(\ulcorner A\urcorner)
              \to \mathrm{Pr}(\ulcorner B\urcorner),
  \quad
  \text{D3} : \vdash \mathrm{Pr}(\ulcorner A\urcorner) \to
              \mathrm{Pr}(\ulcorner \mathrm{Pr}(\ulcorner A\urcorner)\urcorner).
\]
Every premise is a provability statement \emph{about} a code of a
specific formula.  The chain stays at the level ``$F$ is provable'' for
various $F$.  Reflection (D3) is needed to handle the iteration.

Guard's BRA chain operates one level lower: at the level of
\emph{equations} $\mathit{thmT}\;t = u$ for arbitrary Terms $t, u$.
Reflection (D3-analogue) is realised internally and parametrically,
without reference to any specific formula:
\begin{itemize}
  \item Theorem~12 is the parametric ``$D_f$ is the
        encoded image of $f$'' --- one statement, all $f$.
  \item Theorem~13 lifts external equations to encoded form,
        parametric in the equational hypothesis.
  \item Theorem~14 chains substituted-codes through $\mathit{mp}$/
        $\mathit{sb}$/$\mathit{ruleInst}$, never invoking a separate
        ``$\mathrm{Pr}(\ulcorner \mathrm{Pr}(\ulcorner G\urcorner)\urcorner)$''
        reflection axiom.  $\mathit{step5}$ is an equational
        consequence of Theorem~12's defining equations under hypothesis,
        not a provability claim about a specific formula.
\end{itemize}

\subsection{Remarkable consequences}

\begin{enumerate}
  \item \textbf{Avoids meta-reflection.}  No $\mathrm{Pr}(\ulcorner
        \mathrm{Pr}(\ulcorner F\urcorner)\urcorner)$ axiom is invoked.
        The analogous internalisation falls out of $\mathit{thmT}$'s
        defining equations + Comp2/KT/cor compositions.
  \item \textbf{Fully equational.}  Goedel~II reduces to BRA equational
        reasoning + $\mathit{mp}$/$\mathit{axS}$/$\mathit{axK}$/
        $\mathit{axNeg}$/$\mathit{axContrapos}$.  No specialised
        ``provability logic'' is needed.
  \item \textbf{Parametric over $(x : \mathit{Term})$.}  Each step lemma
        works for any Term $x$ (variable, application of cor to a
        variable, anything).  Standard HBL proofs are usually
        meta-iterated over each closed witness.
  \item \textbf{The intermediate Derivs are genuine BRA proofs that do
        not correspond to ``BRA proves formula $F$'' for any single
        $F$.}  This is the structurally interesting feature: the chain
        proves \emph{equational} statements between Terms, with the
        ``provability'' interpretation recovered only at the closure.
\end{enumerate}

The cost is the bookkeeping the BRA edition (and now BRA2) make
explicit: the substituted-code RHS at variable $x$ requires
$\mathit{IsValue}$ certificates (on $j$, on encode-witnesses, on
$t\_term$/$t'\_term$) propagated through to the closure.  These
certificates are the formal counterpart of Guard's silent assumption
that ``$x_$'' for variable $x$ is sensible because $x$ will eventually
be instantiated at a closed numeral.

In short: \emph{Goedel~II in BRA is an equational proof at the code
layer, internalised via primitive recursion, with the
provability-logic story emerging only at closure}.  Guard's 1963
formulation has this character intrinsically; HBL-style proofs do not.
The Agda dev makes the distinction visible.

\section{Methodology}

\subsection{Tree=Term collapse}

The collapse from a separate $\mathit{Tree}$ datatype to
$\mathit{Tree} = \mathit{Term}$ + $\mathit{IsValue}$ predicate was
executed mechanically in May 2026 following BRA2/COLLAPSE-PLAN.md.  All
40-odd files in BRA2/Thm/Parts/, BRA2/Thm12/Parts/, BRA2/Thm12/Param/
were updated by bulk Python transforms to thread $\mathit{IsValue}$
arguments to $\mathit{axKT}$, $\mathit{ktRed}$, $\mathit{constF}_2\mathit{Red}$,
$\mathit{subTDef}$, $\mathit{treeEqRed}$, $\mathit{corOfReify}$,
$\mathit{reifyClosed}$.  The $\mathit{encode}$ contract was strengthened
to expose $\mathit{IsValue}\;y$ alongside the encoded proof tree.

\subsection{Constraints maintained throughout}

\begin{itemize}
  \item \verb|--safe --without-K --exact-split| on every file.
  \item ASCII only.  No Unicode.
  \item No postulates.  No holes.  No \verb|with|-abstraction.  No dot patterns
        (one exception: \verb|.O| / \verb|.a .b| in $\mathit{IsValue}$ pattern matches,
        which are the canonical syntax for matching against an indexed family's
        index).
  \item No new $\mathit{Deriv}$ axioms beyond Guard's primitive list.
\end{itemize}

\subsection{Verification}

\begin{verbatim}
$ agda --safe BRA2/GoedelIIFull.agda
$ echo $?
0
$ ls BRA2/GoedelIIFull.agdai
BRA2/GoedelIIFull.agdai
$ grep -rn "^postulate" BRA2/
$ # (empty)
\end{verbatim}

\section{Discussion}

The two G\"odel theorems in BRA2:
\begin{itemize}
  \item \textbf{G\"odel I.} If the diagonal sentence $G$ is provable, then $\bot$ is provable.
  \item \textbf{G\"odel II.} If the consistency formula $\mathit{Con}$ is provable, then $\bot$ is provable.
\end{itemize}

The intermediate functions $\mathit{constr5}_{\mathrm{final}}$ and the
parametric step lemmas $\mathit{step1\_l}$--$\mathit{step5\_l}$,
$K\_\mathit{part}\_l$, $h\_\mathit{part}\_\mathit{skel}\_l$ are
\emph{first-order} in the sense of recursive-arithmetic definability:
each is a closed combinator built from primitive $\mathit{Fun}_1$ /
$\mathit{Fun}_2$ symbols, and each $\mathit{Deriv}$ derivation is given
by an explicit equational proof in the BRA system.  No external
soundness theorem is invoked.

\subsection{Comparison with prior formalisations}

The only previous machine formalisation of G\"odel's \emph{Second}
Incompleteness Theorem we are aware of is Paulson's Isabelle/HOL
development of hereditarily finite set theory~\cite{paulson2014}.
Other formalisations (O'Connor in Coq, Han in Lean, Carneiro's metamath
work) target G\"odel's \emph{First} theorem only.

Paulson's development uses HF set theory and full G\"odel numbering of
Hilbert-style proofs, with encoded substitution treated as an external
soundness theorem.  The BRA / BRA2 development by contrast follows
Guard~\cite{guard1963}: the ``provability predicate''
$\mathit{thmT}$ is itself a primitive-recursive combinator built from
$\mathit{Rec} + \mathit{Fan} + \mathit{Lift} + \mathit{IfLf}$; encoded
substitution is performed by an explicit primitive recursive function
$\mathit{sub} : \mathit{Fun}_2$; the entire chain stays first-order;
and the proof is equational at the encoded layer rather than
provability-logical at the formula layer
(Section~\ref{sec:internal}).

\begin{thebibliography}{9}

\bibitem{guard1963}
J.~R.~Guard,
\newblock \emph{Lecture notes on recursive arithmetic},
\newblock Applied Mathematics Division, Argonne National Laboratory, 1963.

\bibitem{paulson2014}
L.~C.~Paulson,
\newblock \emph{A machine-assisted proof of G\"odel's incompleteness theorems for the theory of hereditarily finite sets},
\newblock Review of Symbolic Logic, vol.~7 no.~3, 2014, pp.~484--498.

\bibitem{daviespfenning2001}
R.~Davies and F.~Pfenning,
\newblock \emph{A modal analysis of staged computation},
\newblock Journal of the ACM, vol.~48 no.~3, May 2001, pp.~555--604.

\end{thebibliography}

\end{document}
